\documentclass[12pt]{article}
\usepackage{amsmath}
\usepackage{graphicx}

\title{Selective Oxidation of Methanol on Supported Copper Clusters}
\author{Author et al.}
\date{2026}

\begin{document}
\maketitle

\begin{abstract}
Supported copper clusters catalyse the partial oxidation of methanol to formaldehyde.
We compare cluster sizes of four to twelve atoms on a magnesium oxide support using density
functional theory, and we find that six-atom clusters give the highest selectivity among the
sizes studied under the conditions examined here.
\end{abstract}

\section{Introduction}

Formaldehyde production relies on metal oxide catalysts whose active sites remain difficult to
identify at the atomic scale. Supported copper clusters are candidate model catalysts because
their size can be controlled during deposition.

This work examines how cluster size affects the activation barrier for methanol oxidation on a
magnesium oxide support, and it relates the barrier trend to the coordination of surface copper
atoms.

\section{Methods}

All periodic density functional calculations were performed with the Gaussian and plane wave
code. A plane-wave cutoff of 480 eV was used together with a 4x4x4 k-point mesh, and the
Perdew-Burke-Ernzerhof functional described exchange and correlation. Cluster models were cut
from a magnesite surface slab and relaxed until the residual forces fell below a fixed
tolerance.

Adsorption energies were computed with the same settings as the geometry optimizations, and
barriers followed from the climbing-image nudged elastic band method with seven images along
each path.

\section{Discussion}

The six-atom cluster binds the methoxy intermediate more weakly than the larger clusters, which
loosens the carbon-oxygen bond and lowers the barrier for hydrogen abstraction. This
coordination argument is consistent with the d-band picture used for supported metal clusters.

Larger clusters behave like small facets of the bulk metal and lose the size sensitivity of the
barrier, so the discussion of selectivity centres on the under-coordinated corner atoms of the
six-atom cluster.

\section{Results}

Activation barriers for methanol oxidation decrease from 1.05 eV on the four-atom cluster to
0.72 eV on the six-atom cluster, and they rise again to 0.94 eV on the twelve-atom cluster.
The minimum at six atoms is reproduced for both adsorption sites examined.

Selectivity toward formaldehyde follows the same trend, with the six-atom cluster giving the
highest ratio among the sizes studied under the conditions examined here.

\section{Conclusions}

Copper cluster size controls the activation barrier for methanol oxidation on a magnesium oxide
support, and the six-atom cluster is the most selective of the sizes examined. The barrier
trend saturates for clusters larger than eight atoms, which mirrors the loss of size sensitivity
seen in the results.

\end{document}
